\section{Perturbation Distribution}
	Assuming  H II bubbles populate a large volume $V_0$ as a Poisson process.In addition the bubbles' distribution of radii is expressed in an arbitrary distribution $P(R)$.The volume of each bubble $V_b$ is,
	\begin{equation}
	 V_b(R) = \dfrac{4}{3}\pi R^3.
	\end{equation}
	
	Then the ionization fraction at point $\mathbf{r}$ is
	\begin{equation}
		\left<
			x(\textbf{r})
		\right>_P
		=
		1-\exp{
			-\int \dd{R}{
				\dfrac{
					\dd{n_b}}{\dd {R}
					}V_b(R)
				}
			},
		\label{eq:xP}
	\end{equation}
	where ${\dd{n_b}}/{\dd {R}} = n_b(\textbf{r})P(R)$, and $n_b(\textbf{r})$ is \textcolor{blue}{the comoving  bubble number density}. The subscript P means averaging over the Poisson process.
	
	Taking the average of equation \eqref{eq:xP} over volume $V_0$, we can get the averaged ionization fraction in $V_0$
	\begin{equation}
		x_e = 
		1-
		\dfrac{1}{V_0}\int_{V_0}\dd^3{r}~{\exp{-\int \dd{R}{\dfrac{\dd{n_b}}{\dd {R}}V_b(R)}}}
		\label{eq:x_e}
	\end{equation}
	
	Assuming the number density fluctuates as \textcolor{blue}{a bias tracer of the large scale structure} with bubble bias $b(R)$: 
	\begin{equation}
		\dfrac{\dd{n_b}}{\dd{R}}
		=\dfrac{
			\dd{\bar{n}_b}}{\dd{R}
			}
		(1+b(R)
		\delta_W(\textbf{r})),
		\label{eq:nb_fluc}
	\end{equation}
	where $\delta_W(\textbf{r}) = \int \dd^3{r^\prime}{\delta(\textbf{r}^\prime)W_R(\textbf{r}-\textbf{r}^\prime)}$ is the overdensity $\delta$ smoothed with \textcolor{blue}{a top-hat window function}
	\begin{equation}
		W_R(k) = \dfrac{3}{(kR)^3}[\sin(kR)-kR\cos(kR)]
		\label{eq:W_R}
	\end{equation}
	For simplicity, we assume the bubble bias $b(R)$ is a constant b, independent of R.
	
	Substituting with the number density fluctuation  \eqref{eq:nb_fluc} into equation \eqref{eq:x_e} and expanding the exponential term, we obtain:
	\begin{align}
		1-x_e 
		&= 
		\dfrac{1}{V_0}
		\int_{V_0}\dd^3{r}
		~{\exp{
				-\int \dd{R}{
				\dfrac{
					\dd{\bar{n}_b}}{\dd{R}
				}
				(
					1+b\delta_W(\textbf{r})
				)V_b(R)}
			}
		}\notag\\
		&\approx
		\dfrac{1}{V_0}
		\int_{V_0}\dd^3{r}
		~\left[
			1-\int \dd{R}{
				\dfrac{
					\dd{\bar{n}_b}}{\dd{R}
				}
				(
				b\delta_W(\textbf{r})
				)V_b(R)}
		\right]
		\exp{
			-\int\dd{R}{
				\dfrac{
					\dd{\bar{n}_b}
					}{\dd{R}}
				V_b(R)
			}
		}
		\label{eq:expand_x_e}
	\end{align}
	Assuming \textcolor{blue}{the integral of $\delta_W(\textbf{r})$ over a sufficiently large volume $V_0$ approaches 0}, the integral in the square brackets goes to 0.For the additional exponential term, using ${\dd{n_b}}/{\dd {R}} = n_b(\textbf{r})P(R)$, we obtain
	\[
	\exp{
		-\int\dd{R}{
			\dfrac{
				\dd{\bar{n}_b}
			}{\dd{R}}
			V_b(R)
		}
	}=
	\exp{-\bar{n}_b\int\dd{R}{
		V_b(R)P(R)
	}}
	=e^{-\bar{n}_b\left<V_b\right>}
	,
	\]
	where 
	\[
	\left<V_b\right> \equiv \int\dd{R}{
		V_b(R)P(R)
	}.
	\]
	Since the exponential term is independent of the parameter r, it can be taken out of the integral. Therefore the equation \eqref{eq:expand_x_e} can be written as
	\begin{equation}
		1-x_e \approx e^{-\bar{n}_b\left<V_b\right>}
		\label{eq:1-x_e}
	\end{equation}	
	
	Similarly, expand the equation \eqref{eq:xP}, we obtain
	\begin{align}
		\left<
		x(\textbf{r})
		\right>_P
		&=
		1-\exp{
			-\int \dd{R}{
				\dfrac{
					\dd{n_b}}{\dd {R}
				}V_b(R)
			}
		}
		\notag\\
		&=
		1-\exp{
			-\int \dd{R}\left[
				\bar{n}_b P(R)(
				1+b\delta_W(\mathbf{r})
				)
				V_b(R)
			\right]
		}
		\notag\\
		&=
		1-e^{
			-\bar{n}_b\left<
				V_b
			\right>
			}\times\exp{
			-\int \dd{R}\left[
				\bar{n}_b P(R)b\delta_W(\mathbf{r})V_b(R)
			\right]
		}
		\notag\\
		&\approx
		1-e^{
			-\bar{n}_b\left<
				V_b
			\right>
		}\times\left\{
			1-\int \dd{R}\left[
				\bar{n}_b P(R)b\delta_W(\mathbf{r})V_b(R)
			\right]
		\right\}
		\notag\\
		&=
		1-(1-x_e)
		+(1-x_e) b \bar{n}_b
		\times\int \dd{R}\left[
				P(R)\delta_W(\mathbf{r})V_b(R)
			\right]
		\notag\\
		&=
		x_e + (1-x_e) \bar{n}_b b \int\dd{R}{P(R) V_b(R) \delta_W(\mathbf{r})},
	\end{align}
	so the perturbation of the ionized fraction at point $\mathbf{r}$ is
	\begin{align*}
		\left<
		\delta{x(\textbf{r})}
		\right>_P
		=
		\left<
		x(\textbf{r})
		\right>_P
		-x_e
		=(1-x_e) \bar{n}_b b \int\dd{R}{P(R) V_b(R) \delta_W(\mathbf{r})}.
	\end{align*}
	Using the Eq.\eqref{eq:1-x_e}, we obtain $\bar{n}_b = -\dfrac{\ln(1-x_e)} {\left<V_b\right>}$. Substituting $\bar{n}_b$ to $\left<\delta{x(\textbf{r})}\right>_P$,
	we obtain
	\begin{equation}
		\left<
		\delta{x(\textbf{r})}
		\right>_P
		=
		-\dfrac{
			(1-x_e)\ln(1-x_e)
			}{\left<
				V_b
			\right>} b \int\dd{R}{
				P(R) V_b(R) \delta_W(\mathbf{r})
			}
	\end{equation}